%% kapa: exam=1 round=28 year=2017 period=1 form=A subject=economics subject_ko=경제학원론 seq=14 number=54
%% source: Q-Net 감정평가사 기출문제 (cst00309) artlSeq=5199603 file=extracted/1008/5199603/2017년 제28회 감정평가사 1차시험 문제지/1교시 A형.pdf
%% conversion: kordoc 4.24.51 -> tools/gen1.py (auto); review=proofread-sonnet; flags=none
\begin{question}{54}
독점기업의 이윤극대화에 관한 설명으로 옳지 않은 것은? (단, 수요곡선은 우하향하고 생산량은 양(+)이고, 가격차별은 없다.)
\begin{choices}
  \item 이윤극대화 가격은 한계비용보다 높다.
  \item 양(+)의 경제적 이윤을 획득할 수 없는 경우도 있다.
  \item 현재 생산량에서 한계수입이 한계비용보다 높은 상태라면 이윤극대화를 위하여 가격을 인상하여야 한다.
  \item 이윤극대화 가격은 독점 균형거래량에서의 평균수입과 같다.
  \item 이윤극대화는 한계비용과 한계수입이 일치하는 생산수준에서 이루어진다.
\end{choices}
\end{question}
\answer{3}
